\clearpage
\phantomsection
\addcontentsline{toc}{section}{Referencias}
\begin{thebibliography}{99}
\small
\setlength{\itemsep}{0.25em}

\bibitem{qq17} Quan, Q. (2017). \emph{Introduction to Multicopter Design and
Control}. Springer, Singapur. (Caps.~4, 5 y 6: propulsor, marcos y modelo
dinámico.)

\bibitem{qq20} Quan, Q., Dai, X. y Wang, S. (2020). \emph{Multicopter Design
and Control Practice: A Series Experiments Based on MATLAB and Pixhawk}.
Springer, Singapur. ISBN 978-981-15-3137-8.

\bibitem{cs19} Dai, X. CopterSim, archivo \texttt{Init.m}. Repositorio
\texttt{github.com/XunhuaDai/CopterSim}, commit \texttt{0d14487f} (25 de julio
de 2019). Consultado el 4 de octubre de 2026.

\bibitem{hb16} Powers, C., Mellinger, D. y Kumar, V. (2015). Quadrotor
kinematics and dynamics. En K.~P. Valavanis y G.~J. Vachtsevanos (eds.),
\emph{Handbook of Unmanned Aerial Vehicles}, cap.~16, pp.~307--328. Springer.

\bibitem{hb11} Phang, S.~K., Li, K., Chen, B.~M. y Lee, T.~H. (2015).
Systematic design methodology and construction of micro aerial quadrotor
vehicles. En \emph{Handbook of Unmanned Aerial Vehicles}, cap.~11,
pp.~181--206. Springer.

\bibitem{jirinec11} Jiřinec, T. (2011). \emph{Stabilization and Control of
Unmanned Quadcopter}. Tesis de maestría, Czech Technical University in
Prague.

\bibitem{luukkonen11} Luukkonen, T. (2011). \emph{Modelling and Control of
Quadcopter}. Independent research project, Aalto University School of
Science, Espoo.

\bibitem{quitian22} Quitian Niño, L.~C. (2022). \emph{Construcción de bancos
de prueba para la medición experimental de los parámetros del modelo
dinámico de drones -- UAS}. Tesis de maestría, Universidad Nacional de
Colombia, Bogotá.

\bibitem{larotta24} Larotta Mariño, L.~A. (2024). \emph{Control de posición
para un cuadricóptero basado en un sistema de realidad virtual}. Tesis de
maestría, Universidad Nacional de Colombia, Bogotá.

\bibitem{madgwick11} Madgwick, S.~O.~H., Harrison, A.~J.~L. y
Vaidyanathan, R. (2011). Estimation of IMU and MARG orientation using a
gradient descent algorithm. \emph{IEEE International Conference on
Rehabilitation Robotics (ICORR)}.

\bibitem{mahony08} Mahony, R., Hamel, T. y Pflimlin, J.-M. (2008). Nonlinear
complementary filters on the special orthogonal group. \emph{IEEE
Transactions on Automatic Control}, 53(5), 1203--1218.

\bibitem{sabatini06} Sabatini, A.~M. (2006). Quaternion-based extended
Kalman filter for determining orientation by inertial and magnetic sensing.
\emph{IEEE Transactions on Biomedical Engineering}, 53(7), 1346--1356.

\bibitem{leishman14} Leishman, R.~C., Macdonald, J.~C., Beard, R.~W. y
McLain, T.~W. (2014). Quadrotors and accelerometers: State estimation with
an improved dynamic model. \emph{IEEE Control Systems Magazine}, 34(1),
28--41.

\bibitem{skogestad03} Skogestad, S. (2003). Simple analytic rules for model
reduction and PID controller tuning. \emph{Journal of Process Control},
13(4), 291--309.

\bibitem{astrom04} Åström, K.~J. y Hägglund, T. (2004). Revisiting the
Ziegler--Nichols step response method for PID control. \emph{Journal of
Process Control}, 14(6), 635--650.

\bibitem{mellinger11} Mellinger, D. y Kumar, V. (2011). Minimum snap
trajectory generation and control for quadrotors. \emph{IEEE International
Conference on Robotics and Automation (ICRA)}, 2520--2525.

\bibitem{richter13} Richter, C., Bry, A. y Roy, N. (2013). Polynomial
trajectory planning for aggressive quadrotor flight in dense indoor
environments. \emph{International Symposium of Robotics Research (ISRR)}.

\bibitem{cheeseman55} Cheeseman, I.~C. y Bennett, W.~E. (1955). The effect
of the ground on a helicopter rotor in forward flight. \emph{Aeronautical
Research Council, Reports and Memoranda} No.~3021.

\bibitem{mil8785} U.S. Department of Defense (1980). \emph{MIL-F-8785C:
Flying Qualities of Piloted Airplanes}.

\bibitem{milhdbk1797} U.S. Department of Defense (1997). \emph{MIL-HDBK-1797:
Flying Qualities of Piloted Aircraft}.

\bibitem{rbj} Bristow-Johnson, R. \emph{Cookbook formulae for audio EQ
biquad filter coefficients} (``Audio EQ Cookbook''). Documento técnico de
amplia circulación, publicado por la W3C como nota de grupo de trabajo.

\bibitem{ism330} STMicroelectronics. \emph{ISM330DHCX: iNEMO inertial module,
3D accelerometer and 3D gyroscope}. Hoja de datos DS13012.

\bibitem{mmc5983} MEMSIC. \emph{MMC5983MA: 3-axis magnetic sensor}. Hoja de
datos.

\bibitem{bmp180} Bosch Sensortec. \emph{BMP180 digital pressure sensor}. Hoja
de datos BST-BMP180-DS000.

\bibitem{uiuc} Brandt, J.~B., Deters, R.~W., Ananda, G.~K., Dantsker, O.~D. y
Selig, M.~S. \emph{UIUC Propeller Database}. University of Illinois at
Urbana-Champaign.

\bibitem{hairer93} Hairer, E., Nørsett, S.~P. y Wanner, G. (1993).
\emph{Solving Ordinary Differential Equations I: Nonstiff Problems}, 2.ª ed.
Springer.

\bibitem{landau76} Landau, L.~D. y Lifshitz, E.~M. (1976). \emph{Mechanics},
3.ª ed. (Course of Theoretical Physics, vol.~1). Pergamon.

\bibitem{beck09} Beck, A. y Teboulle, M. (2009). A fast iterative
shrinkage-thresholding algorithm for linear inverse problems. \emph{SIAM
Journal on Imaging Sciences}, 2(1), 183--202.

\bibitem{odonoghue15} O'Donoghue, B. y Candès, E. (2015). Adaptive restart
for accelerated gradient schemes. \emph{Foundations of Computational
Mathematics}, 15(3), 715--732.

\bibitem{pannocchia03} Pannocchia, G. y Rawlings, J.~B. (2003). Disturbance
models for offset-free model-predictive control. \emph{AIChE Journal},
49(2), 426--437.

\bibitem{maeder09} Maeder, U., Borrelli, F. y Morari, M. (2009). Linear
offset-free model predictive control. \emph{Automatica}, 45(10),
2214--2222.

\bibitem{li04} Li, W. y Todorov, E. (2004). Iterative linear quadratic
regulator design for nonlinear biological movement systems.
\emph{International Conference on Informatics in Control, Automation and
Robotics (ICINCO)}.

\bibitem{tassa12} Tassa, Y., Erez, T. y Todorov, E. (2012). Synthesis and
stabilization of complex behaviors through online trajectory optimization.
\emph{IEEE/RSJ International Conference on Intelligent Robots and Systems
(IROS)}.

\bibitem{tassa14} Tassa, Y., Mansard, N. y Todorov, E. (2014).
Control-limited differential dynamic programming. \emph{IEEE International
Conference on Robotics and Automation (ICRA)}.

\bibitem{diehl05} Diehl, M., Bock, H.~G. y Schlöder, J.~P. (2005). A
real-time iteration scheme for nonlinear optimization in optimal feedback
control. \emph{SIAM Journal on Control and Optimization}, 43(5),
1714--1736.

\bibitem{torrente21} Torrente, G., Kaufmann, E., Föhn, P. y Scaramuzza, D.
(2021). Data-driven MPC for quadrotors. \emph{IEEE Robotics and Automation
Letters}, 6(2), 3769--3776.

\bibitem{sun22} Sun, S., Romero, A., Foehn, P., Kaufmann, E. y Scaramuzza, D.
(2022). A comparative study of nonlinear MPC and differential-flatness-based
control for quadrotor agile flight. \emph{IEEE Transactions on Robotics},
38(6), 3357--3373.

\bibitem{schulman17} Schulman, J., Wolski, F., Dhariwal, P., Radford, A. y
Klimov, O. (2017). Proximal policy optimization algorithms.
arXiv:1707.06347.

\bibitem{johannink19} Johannink, T., Bahl, S., Nair, A., Luo, J., Kumar, A.,
Loskyll, M., Ojea, J.~A., Solowjow, E. y Levine, S. (2019). Residual
reinforcement learning for robot control. \emph{IEEE International
Conference on Robotics and Automation (ICRA)}.

\bibitem{silver18} Silver, T., Allen, K., Tenenbaum, J. y Kaelbling, L.
(2018). Residual policy learning. arXiv:1812.06298.

\bibitem{tobin17} Tobin, J., Fong, R., Ray, A., Schneider, J., Zaremba, W. y
Abbeel, P. (2017). Domain randomization for transferring deep neural
networks from simulation to the real world. \emph{IEEE/RSJ IROS}.

\bibitem{mckay79} McKay, M.~D., Beckman, R.~J. y Conover, W.~J. (1979). A
comparison of three methods for selecting values of input variables in the
analysis of output from a computer code. \emph{Technometrics}, 21(2),
239--245.

\bibitem{hooke61} Hooke, R. y Jeeves, T.~A. (1961). ``Direct search''
solution of numerical and statistical problems. \emph{Journal of the ACM},
8(2), 212--229.

\bibitem{lawkelton} Law, A.~M. y Kelton, W.~D. \emph{Simulation Modeling and
Analysis}. McGraw-Hill (cap. sobre técnicas de reducción de varianza:
números aleatorios comunes).

\end{thebibliography}
